% =============================================================
%  Module Description — Operating Systems (CS402)
%  DRAFT for a planned module (Semester 4 of the Foundation programme).
%  The code CS402 is a proposal (semester digit + position in the
%  semester); add the term suffix (e.g. -S27) once the term is known and
%  make it equal the course's `code` in src/content/courses/.
%  Compile with: pdflatex cs402-module-description.tex
%  Write "TODO(owner): ..." where information is still missing;
%  the website highlights it.
% =============================================================
\documentclass[11pt,a4paper]{article}

\usepackage[utf8]{inputenc}
\usepackage[T1]{fontenc}
\usepackage{lmodern}
\usepackage[margin=2.2cm]{geometry}
\usepackage{booktabs}
\usepackage{tabularx}
\usepackage{array}
\usepackage{enumitem}
\usepackage{xcolor}
\usepackage{titlesec}
\usepackage{fancyhdr}
\usepackage[hidelinks]{hyperref}

% ---------- Module details (edit these) ----------------------
% \ModuleCode must equal the course's `code` in src/content/courses/.
\newcommand{\ModuleTitle}{Operating Systems}
\newcommand{\ModuleCode}{CS402}
\newcommand{\Programme}{Foundation programme}
\newcommand{\Level}{Semester 4}
\newcommand{\Semester}{Spring}
\newcommand{\Credits}{TODO(owner)}
\newcommand{\ModuleLeader}{Dr. Ali Khudiyev}
\newcommand{\Email}{ali.khudiyev@wosm.academy}
\newcommand{\OfficeHours}{Monday--Friday, 10:00--17:00 (Baku time, UTC+4)}
\newcommand{\Prerequisites}{CS102-F26 Digital Design and Computer Architecture I; CS201 Data Structures and Algorithms I}

% ---------- Styling ------------------------------------------
\definecolor{accent}{RGB}{0,70,130}
\titleformat{\section}{\large\bfseries\color{accent}}{\thesection.}{0.5em}{}
\titlespacing*{\section}{0pt}{1.2em}{0.5em}
\setlist{itemsep=2pt, topsep=4pt}
\setlength{\parindent}{0pt}
\setlength{\parskip}{4pt}

\pagestyle{fancy}
\fancyhf{}
\lhead{\small \ModuleCode\ -- \ModuleTitle}
\rhead{\small Module Description}
\cfoot{\small \thepage}
\renewcommand{\headrulewidth}{0.4pt}

% =============================================================
\begin{document}

\begin{center}
  {\LARGE\bfseries\color{accent} \ModuleTitle}\\[4pt]
  {\large Module Description}
\end{center}
\vspace{6pt}

% ---------- Summary table (not shown on the website: the course page
% shows these details from the macros above) --------------------------
\begin{tabularx}{\textwidth}{>{\bfseries}l X >{\bfseries}l X}
\toprule
Module code   & \ModuleCode    & Credits      & \Credits \\
Programme     & \Programme     & Level        & \Level \\
Semester      & \Semester      & Prerequisites & \Prerequisites \\
Module leader & \ModuleLeader  & Email        & \href{mailto:\Email}{\Email} \\
Office hours  & \multicolumn{3}{l}{\OfficeHours} \\
\bottomrule
\end{tabularx}

% ---------- Sections (shown on the website) ------------------
\section{Module Overview}
An operating system turns bare hardware into a machine that many programs can share
safely and efficiently. This module explains how it does so, around three ideas:
\emph{virtualisation} of the processor and memory, \emph{concurrency}, and
\emph{persistence} of data in file systems.

The module builds on the two Digital Design and Computer Architecture modules, which
explain the hardware the operating system manages: privilege modes, interrupts,
caches and memory. Practical work is in C using the POSIX interface, building on the
C programming from the Data Structures and Algorithms modules.

\section{Aims}
\begin{itemize}
  \item Explain how an operating system virtualises the processor and memory and manages devices and storage.
  \item Teach the principles and pitfalls of concurrent programming.
  \item Show the design trade-offs in scheduling, memory management and file systems.
  \item Give practical experience of systems programming in C.
\end{itemize}

\section{Learning Outcomes}
After completing this course, you will be able to:
\begin{enumerate}[label=LO\arabic*.]
  \item Explain the role and structure of an operating system, including privilege modes, interrupts, traps and system calls.
  \item Explain processes and threads, and evaluate CPU scheduling policies.
  \item Write correct concurrent programs using locks, semaphores and condition variables, and recognise, prevent and handle race conditions and deadlock.
  \item Explain and analyse memory management: address spaces, segmentation, paging, translation lookaside buffers and page replacement.
  \item Explain how file systems and I/O are implemented, including crash consistency.
  \item Write systems programs in C that use the POSIX interface for processes, threads, memory and files.
\end{enumerate}

\section{Syllabus}
\begin{tabularx}{\textwidth}{c X}
\toprule
\textbf{Unit} & \textbf{Topic} \\
\midrule
1  & Introduction: what an operating system does; operating system structures \\
2  & Hardware support: privilege modes, interrupts, traps and system calls \\
3  & Processes and the process API: fork, exec, wait \\
4  & Inter-process communication: pipes, signals and shared memory \\
5  & Threads and POSIX threads \\
6  & CPU scheduling: FCFS, SJF, round robin, MLFQ; multiprocessor scheduling \\
7  & Synchronisation I: race conditions, critical sections, locks \\
8  & Synchronisation II: semaphores, condition variables, monitors; classic problems \\
9  & Deadlock: conditions, prevention, avoidance and detection \\
10 & Memory management: address spaces, segmentation and paging \\
11 & Virtual memory: TLBs, multi-level page tables, page replacement \\
12 & I/O devices, device drivers; hard disks and SSDs \\
13 & File systems: files, directories and their implementation \\
14 & Crash consistency: journaling and log-structured file systems \\
15 & Protection and security \\
16 & Virtual machines and containers \\
17 & Course Review \& Exam \\
\bottomrule
\end{tabularx}

\section{Teaching and Learning Methods}
The module is self-paced. All materials are provided on Moodle.
Each unit has a lecture and exercises. Programming assignments in C let you build
parts of operating-system functionality, such as a shell, a thread-safe data
structure or a memory allocator.

\begin{tabularx}{0.6\textwidth}{X r}
\toprule
\textbf{Activity} & \textbf{Hours} \\
\midrule
Lectures            & 34 \\
Labs / Tutorials    & 46 \\
Independent study   & 145 \\
\midrule
\textbf{Total}      & \textbf{225} \\
\bottomrule
\end{tabularx}

\section{Assessment}
Your knowledge will be assessed throughout the course with quizzes, programming
assignments and the final exam. The exam checks whether you can explain and analyse
operating-system mechanisms and policies and reason about concurrent code. The
programming assignments check whether you can write working systems programs in C.

\begin{tabularx}{\textwidth}{X c c c}
\toprule
\textbf{Component} & \textbf{Weight} & \textbf{Due} & \textbf{LOs} \\
\midrule
Quizzes & 20\% & Throughout & LO1--LO5 \\
Programming assignments & 40\% & Throughout & LO3, LO6 \\
Final exam & 40\% & End of module & LO1--LO5 \\
\midrule
\textbf{Total} & \textbf{100\%} & & \\
\bottomrule
\end{tabularx}

\section{Reading List}
\textbf{Essential}
\begin{itemize}
  \item R.\ H.\ Arpaci-Dusseau, A.\ C.\ Arpaci-Dusseau: \textit{Operating Systems: Three Easy Pieces}, Arpaci-Dusseau Books, 2018. Free online at \url{https://pages.cs.wisc.edu/~remzi/OSTEP/}.
\end{itemize}
\textbf{Recommended}
\begin{itemize}
  \item A.\ Silberschatz, P.\ B.\ Galvin, G.\ Gagne: \textit{Operating System Concepts}, 10th ed., Wiley, 2018.
  \item A.\ S.\ Tanenbaum, H.\ Bos: \textit{Modern Operating Systems}, 4th ed., Pearson, 2014.
  \item R.\ E.\ Bryant, D.\ R.\ O'Hallaron: \textit{Computer Systems: A Programmer's Perspective}, 3rd ed., Pearson, 2016.
  \item B.\ W.\ Kernighan, D.\ M.\ Ritchie: \textit{The C Programming Language}, 2nd ed., Prentice Hall, 1988.
\end{itemize}

\section{Materials and Licence}
Course materials are published under the
\href{https://creativecommons.org/licenses/by-sa/4.0/}{Creative Commons Attribution-ShareAlike 4.0}
licence unless stated otherwise. See \href{https://wosm.academy/licence/}{wosm.academy/licence} for details.

\section{Policies}
\begin{itemize}
  \item \textbf{Attendance:} Attendance is not mandatory. The module is self-paced, so you can study the materials when it suits you.
  \item \textbf{Late submissions:} Work submitted after the deadline is penalised.
  \item \textbf{Academic integrity:} All submitted work must be your own.
        See the \href{https://wosm.academy/code-of-conduct/}{Code of Conduct}.
\end{itemize}

\end{document}
